\documentclass[12pt]{article}
\usepackage{amsmath,amssymb,amsfonts,amsthm}
\usepackage[margin=1in]{geometry}

\begin{document}

\begin{center}
{\Large\bfseries Chapter 9, Problem 26}\\[6pt]
Byron \& Fuller, \textit{Mathematics of Classical and Quantum Physics}
\end{center}

\bigskip

\textbf{Problem.} Suppose that $\{\phi_n\}$ and $\{\psi_m\}$ are sequences of step functions and that $\phi_n \to f$ and $\psi_m \to g$. Show that for any step functions $\phi$ and $\psi$:
\[
(V\phi, \psi) = \langle\phi, U\psi\rangle, \qquad (V\phi, V\psi) = \langle\phi, \psi\rangle, \qquad (U\phi, U\psi) = \langle\phi, \psi\rangle.
\]
Then show that $\{V\phi_n\}$ and $\{U\psi_m\}$ converge to limits $Vf$ and $Ug$, and $(Vf, g) = \langle f, Ug\rangle$. This completes the proof of Theorem~9.11.

\bigskip
\textbf{Solution.}

This problem concerns extending the Fourier transform operators $U$ and $V$ from step functions to all of $L_2$, thereby completing the proof that the Fourier transform defines a unitary transformation on Hilbert space (Theorem~9.11).

\subsection*{Step 1: Properties on step functions}

For step functions, $U$ and $V$ are defined by the Fourier transform and its inverse. The three identities are consequences of the Parseval--Plancherel theorem for step functions:

\textbf{(i) $(V\phi, V\psi) = \langle\phi, \psi\rangle$:}

For step functions $\phi, \psi$, $V\phi$ and $V\psi$ are their Fourier transforms (or inverse transforms). By Parseval's theorem (which holds for step functions by direct computation):
\[
\int_{-\infty}^{\infty} (V\phi)(x)\,\overline{(V\psi)(x)}\,dx = \int_{-\infty}^{\infty} \phi(y)\,\overline{\psi(y)}\,dy.
\]
Thus $(V\phi, V\psi) = \langle\phi, \psi\rangle$. \hfill $\square$

\textbf{(ii) $(U\phi, U\psi) = \langle\phi, \psi\rangle$:} Same argument with $U$ in place of $V$. \hfill $\square$

\textbf{(iii) $(V\phi, \psi) = \langle\phi, U\psi\rangle$:}

This is the adjoint relation. By definition of $V$ and $U$ as Fourier transform/inverse:
\begin{align}
(V\phi, \psi) &= \int (V\phi)(x)\,\overline{\psi(x)}\,dx = \int\left[\frac{1}{\sqrt{2\pi}}\int \phi(y)\,e^{-ixy}\,dy\right]\overline{\psi(x)}\,dx \\
&= \int \phi(y)\left[\frac{1}{\sqrt{2\pi}}\int \overline{\psi(x)}\,e^{-ixy}\,dx\right]dy = \int\phi(y)\,\overline{(U\psi)(y)}\,dy = \langle\phi, U\psi\rangle,
\end{align}
where we exchanged the order of integration (justified for step functions) and used $(U\psi)(y) = \frac{1}{\sqrt{2\pi}}\int \psi(x)\,e^{ixy}\,dx$. \hfill $\square$

\subsection*{Step 2: Extension to $L_2$}

Since $\phi_n \to f$ in $L_2$, and $V$ is an isometry on step functions ($\|V\phi\| = \|\phi\|$):
\[
\|V\phi_n - V\phi_m\| = \|\phi_n - \phi_m\| \to 0 \quad \text{as } n, m \to \infty.
\]
So $\{V\phi_n\}$ is a Cauchy sequence in $L_2$. Since $L_2$ is complete, $V\phi_n \to Vf$ for some $Vf \in L_2$. This defines $Vf$ for all $f \in L_2$.

Similarly, $U\psi_m \to Ug$ for some $Ug \in L_2$.

\subsection*{Step 3: The identity $(Vf, g) = \langle f, Ug\rangle$}

For step functions, $(V\phi_n, \psi_m) = \langle\phi_n, U\psi_m\rangle$. Taking limits:
\begin{itemize}
\item Left side: $V\phi_n \to Vf$ and $\psi_m \to g$ in $L_2$, so $(V\phi_n, \psi_m) \to (Vf, g)$ (inner product is continuous).
\item Right side: $\phi_n \to f$ and $U\psi_m \to Ug$, so $\langle\phi_n, U\psi_m\rangle \to \langle f, Ug\rangle$.
\end{itemize}

Therefore:
\[
\boxed{(Vf, g) = \langle f, Ug\rangle \quad \text{for all } f, g \in L_2}.
\]

Similarly, from $(V\phi_n, V\phi_m) = \langle\phi_n, \phi_m\rangle$, taking limits gives $(Vf, Vg) = \langle f, g\rangle$, proving $V$ is an isometry on all of $L_2$. The same holds for $U$. Since $V^\dagger = U$ (from $(Vf,g) = (f, Ug)$) and both are isometries, $U$ and $V$ are unitary transformations. This completes the proof of Theorem~9.11. \hfill $\blacksquare$

\end{document}
